\icode{N18J-MC-SI-L2-261026-004}
\nceeID[6]{2018北京}{20}
根据高中所学知识可知，做自由落体运动的小球，将落在正下方位置。但实际上，赤道上方\qty{200}{m}处无初速下落的小球将落在正下方位置偏东约\qty{6}{cm}处。这一现象可解释为，除重力外，由于地球自转，下落过程小球还受到一个水平向东的“力”，该“力”与竖直方向的速度大小成正比。现将小球从赤道地面竖直上抛，考虑对称性，上升过程该“力”水平向西，则小球\paren[]
\begin{choices}
	\item 到最高点时，水平方向的加速度和速度均为零
	\item 到最高点时，水平方向的加速度和速度均不为零
	\item 落地点在抛出点东侧
	\item 落地点在抛出点西侧
\end{choices}
\begin{answer}
	D\qquad 提示：考查同学们接受理解信息的能力
\end{answer}